% GENERATED FILE. Edit canonical lessons, then run npm run generate:books.
% canonical-source-hash: fnv1a64:c9c9cd8a5e0e4503
% canonical-lessons: PT-C429-berco, PT-C429-fronha, PT-C429-secador, PT-C429-cabide, PT-C429-roupao

\chapter{No quarto}
\label{ch:pt-no-quarto}
\hlchaptermodality{429}

\begin{chapteropening}
\textbf{By the end of this chapter:} \emph{I can ask for things in my room: a cot, a pillowcase, a hair dryer, a hanger and a bathrobe.}

Complete the last lesson of chapter 429: I can ask for things in my room: a cot, a pillowcase, a hair dryer, a hanger and a bathrobe.
\end{chapteropening}

\section[o berço]{o berço — the cot, the cradle}

\label{lesson:PT-C429-berco}



\begin{quote}
Before the new one: say the Portuguese for to put up, then the Portuguese for the doorman.
\end{quote}

\subsection*{You'll want to know: o berço}
\textbf{o berço} — \textquotedblleft{}the cot, the cradle\textquotedblright{}.

\begin{quote}
Pedimos um berço para o bebé.
\end{quote}

\begin{quote}
\emph{We asked for a cot for the baby.}
\end{quote}

\emph{Nascer em berço de ouro} — to be born with a silver spoon in your mouth.

\begin{cousinweb}
\emph{berço} comes from Vulgar Latin \emph{berciolum} 'cradle' (EtymDB), probably from a Gaulish word; the details are not settled.
\end{cousinweb}

\begin{grammarlens}[title={What you've built}]
One more word for telling, explaining and describing at B1.
\end{grammarlens}

\subsection*{Your turn}
\begin{itemize}
\raggedright
  \item \emph{Say it:} \emph{o berço}
  \item \emph{Say it:} \emph{o berço}, once more
  \item \emph{Say it:} a sentence of your own with \emph{o berço}
  \item \emph{Recall:} say the Portuguese for to put up, then the Portuguese for the doorman, then say \emph{o berço} again
\end{itemize}

\begin{tcolorbox}[breakable,colback=teal!4,colframe=teal!35!black,title={Before you move on}]
What does o berço mean? (\textquotedblleft{}The cot, the cradle\textquotedblright{}.) Say it once more.
\end{tcolorbox}
\section[a fronha]{a fronha — the pillowcase}

\label{lesson:PT-C429-fronha}



\begin{quote}
Before the new one: say the Portuguese for the doorman, then the Portuguese for the cot.
\end{quote}

\subsection*{You'll want to know: a fronha}
\textbf{a fronha} — \textquotedblleft{}the pillowcase\textquotedblright{}.

\begin{quote}
A fronha estava suja e pedi outra.
\end{quote}

\begin{quote}
\emph{The pillowcase was dirty and I asked for another.}
\end{quote}

Sheets are \emph{os lençóis}.

\begin{cousinweb}
The origin of \emph{fronha} is not settled, and EtymDB has no source for it.
\end{cousinweb}

\begin{grammarlens}[title={What you've built}]
One more word for telling, explaining and describing at B1.
\end{grammarlens}

\subsection*{Your turn}
\begin{itemize}
\raggedright
  \item \emph{Say it:} \emph{a fronha}
  \item \emph{Say it:} \emph{a fronha}, once more
  \item \emph{Say it:} a sentence of your own with \emph{a fronha}
  \item \emph{Recall:} say the Portuguese for the doorman, then the Portuguese for the cot, then say \emph{a fronha} again
\end{itemize}

\begin{tcolorbox}[breakable,colback=teal!4,colframe=teal!35!black,title={Before you move on}]
What does a fronha mean? (\textquotedblleft{}The pillowcase\textquotedblright{}.) Say it once more.
\end{tcolorbox}
\section[o secador]{o secador — the hair dryer; the dryer}

\label{lesson:PT-C429-secador}



\begin{quote}
Before the new one: say the Portuguese for the cot, then the Portuguese for the pillowcase.
\end{quote}

\subsection*{You'll want to know: o secador}
\textbf{o secador} — \textquotedblleft{}the hair dryer; the dryer\textquotedblright{}.

\begin{quote}
O secador da casa de banho não funciona.
\end{quote}

\begin{quote}
\emph{The hair dryer in the bathroom does not work.}
\end{quote}

\emph{O secador de cabelo} in full; \emph{a máquina de secar} dries clothes.

\begin{cousinweb}
\emph{secador} is built on \emph{secar} 'to dry', which is inherited from Latin \emph{siccāre}, from \emph{siccus} 'dry' (Lewis \& Short).
\end{cousinweb}

\begin{grammarlens}[title={What you've built}]
One more word for telling, explaining and describing at B1.
\end{grammarlens}

\subsection*{Your turn}
\begin{itemize}
\raggedright
  \item \emph{Say it:} \emph{o secador}
  \item \emph{Say it:} \emph{o secador}, once more
  \item \emph{Say it:} a sentence of your own with \emph{o secador}
  \item \emph{Recall:} say the Portuguese for the cot, then the Portuguese for the pillowcase, then say \emph{o secador} again
\end{itemize}

\begin{tcolorbox}[breakable,colback=teal!4,colframe=teal!35!black,title={Before you move on}]
What does o secador mean? (\textquotedblleft{}The hair dryer; the dryer\textquotedblright{}.) Say it once more.
\end{tcolorbox}
\section[o cabide]{o cabide — the coat hanger; the coat rack}

\label{lesson:PT-C429-cabide}



\begin{quote}
Before the new one: say the Portuguese for the pillowcase, then the Portuguese for the hair dryer.
\end{quote}

\subsection*{You'll want to know: o cabide}
\textbf{o cabide} — \textquotedblleft{}the coat hanger; the coat rack\textquotedblright{}.

\begin{quote}
No armário só há dois cabides.
\end{quote}

\begin{quote}
\emph{There are only two hangers in the wardrobe.}
\end{quote}

In Brazil: \emph{o cabide} too.

\begin{cousinweb}
The origin of \emph{cabide} is not settled; it is sometimes traced to Arabic, and EtymDB has no entry for it.
\end{cousinweb}

\begin{grammarlens}[title={What you've built}]
One more word for telling, explaining and describing at B1.
\end{grammarlens}

\subsection*{Your turn}
\begin{itemize}
\raggedright
  \item \emph{Say it:} \emph{o cabide}
  \item \emph{Say it:} \emph{o cabide}, once more
  \item \emph{Say it:} a sentence of your own with \emph{o cabide}
  \item \emph{Recall:} say the Portuguese for the pillowcase, then the Portuguese for the hair dryer, then say \emph{o cabide} again
\end{itemize}

\begin{tcolorbox}[breakable,colback=teal!4,colframe=teal!35!black,title={Before you move on}]
What does o cabide mean? (\textquotedblleft{}The coat hanger; the coat rack\textquotedblright{}.) Say it once more.
\end{tcolorbox}
\section[o roupão]{o roupão — the bathrobe, the dressing gown}

\label{lesson:PT-C429-roupao}



\begin{quote}
Before the new one: say the Portuguese for the hair dryer, then the Portuguese for the coat hanger.
\end{quote}

\subsection*{You'll want to know: o roupão}
\textbf{o roupão} — \textquotedblleft{}the bathrobe, the dressing gown\textquotedblright{}.

\begin{quote}
No hotel há um roupão branco para cada hóspede.
\end{quote}

\begin{quote}
\emph{At the hotel there is a white bathrobe for each guest.}
\end{quote}

Built on \emph{a roupa} with \emph{-ão}.

\begin{cousinweb}
\emph{roupão} is built on \emph{roupa} 'clothes' with the big-size ending \emph{-ão}; \emph{roupa} comes from Gothic \emph{raupa} (EtymDB).
\end{cousinweb}

\begin{grammarlens}[title={What you've built}]
One more word for telling, explaining and describing at B1.
\end{grammarlens}

\subsection*{Your turn}
\begin{itemize}
\raggedright
  \item \emph{Say it:} \emph{o roupão}
  \item \emph{Say it:} \emph{o roupão}, once more
  \item \emph{Say it:} a sentence of your own with \emph{o roupão}
  \item \emph{Recall:} say the Portuguese for the hair dryer, then the Portuguese for the coat hanger, then say \emph{o roupão} again
\end{itemize}

\begin{tcolorbox}[breakable,colback=teal!4,colframe=teal!35!black,title={Before you move on}]
What does o roupão mean? (\textquotedblleft{}The bathrobe, the dressing gown\textquotedblright{}.) Say it once more.
\end{tcolorbox}
